% Part: first-order-logic
% Chapter: sequent-calculus
% Section: proving-things

\documentclass[../../../include/open-logic-section]{subfiles}

\begin{document}

\iftag{FOL}
      {\olfileid[fr]{fol}{seq}{pro}}
      {\olfileid[fr]{pl}{seq}{pro}}

\olsection{Exemples de \usetoken{p}{derivation}}

\begin{ex}
Construire une !!{derivation} dans $\Log{LK}$ du séquent
$!A \land !B \Sequent !A$.

Commençons par écrire le séquent final recherché au bas de la
!!{derivation}.
\begin{prooftree}
\AxiomC{}
\UnaryInf$!A\land !B \fCenter !A$
\end{prooftree}
Cherchons ensuite une inférence dont le séquent inférieur pourrait
avoir cette forme. Une règle structurelle conviendrait peut-être,
mais il est judicieux de commencer par chercher une règle logique.
Le seul connecteur logique qui figure dans le séquent inférieur
est $\land$ : nous cherchons donc une règle pour $\land$.
Comme ce symbole figure dans l'antécédent, examinons la règle
\LeftR{\land}.
\begin{prooftree}
\AxiomC{}
\RightLabel{\LeftR{\land}}
\UnaryInf$!A\land !B \fCenter !A$
\end{prooftree}
Deux séquents pourraient servir de prémisse à cette inférence
\LeftR{\land} : $!A \Sequent !A$ ou $!B \Sequent !A$.
Le premier est un séquent initial, ce qui nous convient, tandis que
le second n'est pas dérivable en général. Inscrivons donc la prémisse :
\begin{prooftree}
\Axiom$!A \fCenter !A$
\RightLabel{\LeftR{\land}}
\UnaryInf$!A\land !B \fCenter !A$
\end{prooftree}
Nous avons maintenant une !!{derivation} correcte, dans $\Log{LK}$,
du séquent $!A \land !B \Sequent !A$.
\end{ex}

\begin{ex}
Construire une !!{derivation} dans $\Log{LK}$ du séquent
$\lnot !A \lor !B \Sequent !A \lif !B$.

Commençons par écrire le séquent final recherché au bas de la
!!{derivation}.
\begin{prooftree}
\AxiomC{}
\UnaryInf$\lnot !A \lor !B \fCenter !A \lif !B$
\end{prooftree}
Pour trouver une règle logique susceptible de donner ce séquent final,
examinons les connecteurs qui y figurent : $\lnot$, $\lor$ et $\lif$.
Pour l'instant, seuls $\lor$ et $\lif$ nous intéressent, car ce sont
les !!{main operator}s des !!{sentence}s du séquent final.
Le symbole $\lnot$ se trouve dans la portée d'un autre connecteur :
nous nous en occuperons donc plus tard. Les règles logiques possibles
pour la dernière inférence sont ainsi \LeftR{\lor} et \RightR{\lif}.
L'une ou l'autre conviendrait ; choisissons \RightR{\lif}, ne serait-ce
que pour différer la séparation en deux branches. D'après le schéma
de cette règle et le séquent inférieur recherché, l'inférence doit
avoir la forme suivante :
\begin{prooftree}
\AxiomC{}
\UnaryInf$ !A, \lnot !A \lor !B \fCenter !B $
\RightLabel{\RightR{\lif}} \UnaryInf$ \lnot !A \lor !B \fCenter !A \lif !B $
\end{prooftree}

En déplaçant $\lnot !A \lor !B$ à l'extrémité gauche de l'antécédent,
nous pouvons appliquer \LeftR{\lor}. Le schéma de cette règle nous
conduit aux deux prémisses suivantes :\footnote{Note de l'édition française :
dans cet exemple, quatre étapes échangent deux énoncés de l'antécédent.
Leurs étiquettes $\RightR{\Exchange}$ dans la source sont corrigées
en $\LeftR{\Exchange}$.}
\begin{prooftree}
\AxiomC{}
\UnaryInf$\lnot !A, !A \fCenter !B$
\AxiomC{}
\UnaryInf$!B, !A \fCenter !B$
\RightLabel{\LeftR{\lor}}
\BinaryInf$ \lnot !A \lor !B, !A \fCenter !B $
\RightLabel{\LeftR{\Exchange}}
\UnaryInf$ !A, \lnot !A \lor !B \fCenter !B $
\RightLabel{\RightR{\lif}}
\UnaryInf$ \lnot !A \lor !B \fCenter !A \lif !B $
\end{prooftree}
Rappelons que nous cherchons à remonter jusqu'à des séquents initiaux.
Nous en sommes proches : un affaiblissement et un échange suffisent
à achever la branche de droite à partir d'un séquent initial :
\begin{prooftree}
\AxiomC{}
\UnaryInf$\lnot !A, !A \fCenter !B$
\Axiom$!B \fCenter !B$
\RightLabel{\LeftR{\Weakening}}
\UnaryInf$!A, !B \fCenter !B$
\RightLabel{\LeftR{\Exchange}}
\UnaryInf$!B, !A \fCenter !B$
\RightLabel{\LeftR{\lor}}
\BinaryInf$\lnot !A \lor !B, !A \fCenter !B $
\RightLabel{\LeftR{\Exchange}}
\UnaryInf$ !A, \lnot !A \lor !B \fCenter !B $
\RightLabel{\RightR{\lif}}
\UnaryInf$ \lnot !A \lor !B \fCenter !A \lif !B $
\end{prooftree}

Examinons maintenant la branche de gauche. Le seul connecteur logique
qui figure dans ses !!{sentence}s est le symbole $\lnot$, dans
l'antécédent. Nous cherchons donc une instance de \LeftR{\lnot}.
\begin{prooftree}
\AxiomC{}
\UnaryInf$ !A \fCenter !B, !A$
\RightLabel{\LeftR{\lnot}}
\UnaryInf$\lnot !A, !A \fCenter !B$
\Axiom$!B \fCenter !B$
\RightLabel{\LeftR{\Weakening}}
\UnaryInf$!A, !B \fCenter !B$
\RightLabel{\LeftR{\Exchange}}
\UnaryInf$!B, !A \fCenter !B$
\RightLabel{\LeftR{\lor}}
\BinaryInf$\lnot !A \lor !B, !A \fCenter !B $
\RightLabel{\LeftR{\Exchange}}
\UnaryInf$ !A, \lnot !A \lor !B \fCenter !B $
\RightLabel{\RightR{\lif}}
\UnaryInf$ \lnot !A \lor !B \fCenter !A \lif !B $
\end{prooftree}
Comme pour la branche de droite, un affaiblissement et un échange
suffisent à achever cette branche.
\begin{prooftree}
\Axiom$!A \fCenter !A$
\RightLabel{\RightR{\Weakening}}
\UnaryInf$ !A \fCenter !A, !B$
\RightLabel{\RightR{\Exchange}}
\UnaryInf$ !A \fCenter !B, !A$
\RightLabel{\LeftR{\lnot}}
\UnaryInf$\lnot !A, !A \fCenter !B$
\Axiom$!B \fCenter !B$
\RightLabel{\LeftR{\Weakening}}
\UnaryInf$!A, !B \fCenter !B$
\RightLabel{\LeftR{\Exchange}}
\UnaryInf$!B, !A \fCenter !B$
\RightLabel{\LeftR{\lor}}
\BinaryInf$\lnot !A \lor !B, !A \fCenter !B $
\RightLabel{\LeftR{\Exchange}}
\UnaryInf$ !A, \lnot !A \lor !B \fCenter !B $
\RightLabel{\RightR{\lif}}
\UnaryInf$ \lnot !A \lor !B \fCenter !A \lif !B $
\end{prooftree}
\end{ex}

\begin{ex}
Construire une !!{derivation} dans $\Log{LK}$ du séquent
$\lnot !A \lor \lnot !B \Sequent \lnot (!A \land !B)$.

Comme précédemment, commençons par écrire le séquent final recherché
au bas de l'arbre.
\begin{prooftree}
\AxiomC{}
\UnaryInf$ \lnot !A \lor \lnot !B \fCenter \lnot (!A \land !B) $
\end{prooftree}
Les connecteurs principaux des !!{sentence}s du séquent final sont
$\lor$ et $\lnot$. Les règles \LeftR{\lor} et \RightR{\lnot}
conviendraient toutes deux ; commençons par \RightR{\lnot} pour
différer la séparation en deux branches :
\begin{prooftree}
\AxiomC{}
\UnaryInf$!A \land !B, \lnot !A \lor \lnot !B \fCenter $
\RightLabel{\RightR{\lnot}}
\UnaryInf$\lnot !A \lor \lnot !B \fCenter \lnot (!A \land !B)$
\end{prooftree}
Nous avons maintenant le choix entre \LeftR{\land} et \LeftR{\lor}.
Essayons \LeftR{\land}. Nous pouvons partir du séquent
$!A, \lnot !A \lor \lnot !B \Sequent \quad$ ou du séquent
$!B, \lnot !A \lor \lnot !B \Sequent \quad$.
La situation étant symétrique en $!A$ et $!B$, choisissons le premier :\footnote{Note
de l'édition française : les deux prémisses proposées dans la prose source
omettent la négation devant $!B$ dans la disjonction. Elle est rétablie
ici, conformément aux arbres de la source.}
\begin{prooftree}
\AxiomC{}
\UnaryInf$!A, \lnot !A \lor \lnot !B \fCenter $
\RightLabel{\LeftR{\land}}
\UnaryInf$!A \land !B, \lnot !A \lor \lnot !B \fCenter $
\RightLabel{\RightR{\lnot}}
\UnaryInf$\lnot !A \lor \lnot !B \fCenter \lnot (!A \land !B)$
\end{prooftree}
En poursuivant la construction de la !!{derivation}, nous rencontrons
une difficulté :
\begin{prooftree}
\Axiom$!A \fCenter !A$
\RightLabel{\LeftR{\lnot}}
\UnaryInf$ \lnot !A, !A \fCenter$
\AxiomC{}
\RightLabel{?}
\UnaryInf$!A \fCenter !B$
\RightLabel{\LeftR{\lnot}}
\UnaryInf$ \lnot !B, !A \fCenter$
\RightLabel{\LeftR{\lor}}
\BinaryInf$\lnot !A \lor \lnot !B, !A \fCenter $
\RightLabel{\LeftR{\Exchange}}
\UnaryInf$!A, \lnot !A \lor \lnot !B \fCenter $
\RightLabel{\LeftR{\land}}
\UnaryInf$!A \land !B, \lnot !A \lor \lnot !B \fCenter $
\RightLabel{\RightR{\lnot}}
\UnaryInf$\lnot !A \lor \lnot !B \fCenter \lnot (!A \land !B)$
\end{prooftree}
En général, le séquent au sommet de la branche de droite ne peut être
réduit davantage ni obtenu à partir d'un séquent initial par des
inférences structurelles. Cette voie ne convient donc pas : nous avons
appliqué \LeftR{\land} trop tôt. Revenons à l'étape précédente et
appliquons plutôt \LeftR{\lor} :\footnote{Note de l'édition française :
l'échec décrit concerne le schéma général. Certaines instances
particulières de $!A \Sequent !B$, notamment lorsque $!A$ et $!B$
sont identiques, sont dérivables.}
\begin{prooftree}
\AxiomC{}
\UnaryInf$\lnot !A, !A \land !B \fCenter $

\AxiomC{}
\UnaryInf$\lnot !B, !A \land !B \fCenter $

\RightLabel{\LeftR{\lor}}
\BinaryInf$\lnot !A \lor \lnot !B, !A \land !B \fCenter $
\RightLabel{\LeftR{\Exchange}}
\UnaryInf$!A \land !B, \lnot !A \lor \lnot !B \fCenter $
\RightLabel{\RightR{\lnot}}
\UnaryInf$\lnot !A \lor \lnot !B \fCenter \lnot (!A \land !B)$
\end{prooftree}
En achevant chaque branche comme précédemment, nous obtenons :
\begin{prooftree}
\Axiom$ !A \fCenter!A$
\RightLabel{\LeftR{\land}}
\UnaryInf$!A \land !B \fCenter !A$
\RightLabel{\LeftR{\lnot}}
\UnaryInf$\lnot !A, !A \land !B \fCenter $

\Axiom$ !B \fCenter !B$
\RightLabel{\LeftR{\land}}
\UnaryInf$!A \land !B \fCenter !B$
\RightLabel{\LeftR{\lnot}}
\UnaryInf$\lnot !B, !A \land !B \fCenter $

\RightLabel{\LeftR{\lor}}
\BinaryInf$\lnot !A \lor \lnot !B, !A \land !B \fCenter $
\RightLabel{\LeftR{\Exchange}}
\UnaryInf$!A \land !B, \lnot !A \lor \lnot !B \fCenter $
\RightLabel{\RightR{\lnot}}
\UnaryInf$\lnot !A \lor \lnot !B \fCenter \lnot (!A \land !B)$
\end{prooftree}
(On aurait aussi pu placer les règles pour $\land$ plus bas que les
règles pour $\lnot$ dans ces étapes et obtenir une !!{derivation}
correcte.)
\end{ex}

\begin{ex}
Nous n'avons pas encore utilisé la contraction, mais elle est parfois
nécessaire. En voici un exemple. Supposons que nous voulions démontrer
$\quad \Sequent !A \lor \lnot !A$. En appliquant
$\RightR{\lor}$ à rebours, nous obtenons l'une des deux tentatives
de !!{derivation} suivantes :
\begin{prooftree}
\AxiomC{}
\UnaryInf$ \fCenter !A$
\RightLabel{\RightR{\lor}}
\UnaryInf$ \fCenter !A \lor \lnot !A$
\DisplayProof\qquad\bottomAlignProof
\AxiomC{}
\UnaryInf$!A \fCenter $
\RightLabel{\RightR{\lnot}}
\UnaryInf$ \fCenter \lnot !A$
\RightLabel{\RightR{\lor}}
\UnaryInf$ \fCenter !A \lor \lnot !A$
\end{prooftree}
Aucune de ces deux tentatives n'aboutit, en général, à un séquent
initial. L'idée est d'utiliser la contraction, qui réunit deux
occurrences d'!!a{sentence} en une seule. Lors de la recherche d'une
preuve, c'est-à-dire en allant du bas vers le haut, nous pouvons donc
conserver une occurrence de $!A \lor \lnot !A$ dans la prémisse,
par exemple :
\begin{prooftree}
\AxiomC{}
\UnaryInf$ \fCenter !A \lor \lnot !A, !A$
\RightLabel{\RightR{\lor}}
\UnaryInf$ \fCenter !A \lor \lnot !A, !A \lor \lnot !A$
\RightLabel{\RightR{\Contraction}}
\UnaryInf$ \fCenter !A \lor \lnot !A$
\end{prooftree}
Nous pouvons maintenant appliquer $\RightR{\lor}$ une seconde fois,
en faisant aussi apparaître $\lnot !A$, ce qui donne une
!!{derivation} complète :
\begin{prooftree}
\Axiom$!A \fCenter !A$
\RightLabel{\RightR{\lnot}}
\UnaryInf$\fCenter !A, \lnot !A$
\RightLabel{\RightR{\lor}}
\UnaryInf$\fCenter !A, !A \lor \lnot !A$
\RightLabel{\RightR{\Exchange}}
\UnaryInf$ \fCenter !A \lor \lnot !A, !A$
\RightLabel{\RightR{\lor}}
\UnaryInf$ \fCenter !A \lor \lnot !A, !A \lor \lnot !A$
\RightLabel{\RightR{\Contraction}}
\UnaryInf$ \fCenter !A \lor \lnot !A$
\end{prooftree}
\end{ex}

\begin{prob}
Construire des !!{derivation}s des séquents suivants :
\begin{enumerate}
\item $!A \land (!B \land !C) \Sequent (!A \land !B) \land !C$.
\item $!A \lor (!B \lor !C) \Sequent (!A \lor !B) \lor !C$.
\item $!A \lif (!B \lif !C) \Sequent !B \lif (!A \lif !C)$.
\item $!A \Sequent \lnot\lnot !A$.
\end{enumerate}
\end{prob}

\begin{prob}
Construire des !!{derivation}s des séquents suivants :
\begin{enumerate}
\item $(!A \lor !B) \lif !C \Sequent !A \lif !C$.
\item $(!A \lif !C) \land (!B \lif !C) \Sequent (!A \lor !B) \lif !C$.
\item $\Sequent \lnot(!A \land \lnot !A)$.
\item $!B \lif !A \Sequent \lnot !A \lif \lnot !B$.
\item $\Sequent (!A \lif \lnot !A) \lif \lnot !A$.
\item $\Sequent \lnot(!A \lif !B) \lif \lnot !B$.
\item $!A \lif !C \Sequent \lnot (!A \land \lnot !C)$.
\item $!A \land \lnot !C \Sequent \lnot (!A \lif !C)$.
\item $!A \lor !B, \lnot !B \Sequent !A$.
\item $\lnot !A \lor \lnot !B \Sequent \lnot(!A \land !B)$.
\item $\Sequent (\lnot !A \land \lnot !B) \lif\lnot(!A \lor !B)$.
\item $\Sequent \lnot(!A \lor !B) \lif (\lnot !A \land \lnot !B)$.
\end{enumerate}
\end{prob}

\begin{prob}
Construire des !!{derivation}s des séquents suivants :
\begin{enumerate}
\item $\lnot(!A \lif !B) \Sequent !A$.
\item $\lnot(!A \land !B) \Sequent \lnot !A \lor \lnot !B$.
\item $!A \lif !B \Sequent \lnot !A \lor !B$.
\item $\Sequent \lnot \lnot !A \lif !A$.
\item $!A \lif !B, \lnot !A \lif !B \Sequent !B$.
\item $(!A \land !B) \lif !C \Sequent (!A \lif !C) \lor (!B \lif !C)$.
\item $(!A \lif !B) \lif !A \Sequent !A$.
\item $\Sequent (!A \lif !B) \lor (!B \lif !C)$.
\end{enumerate}
(Tous requièrent la règle $\RightR{\Contraction}$.)
\end{prob}
\end{document}
